%
%
% This is an example LaTeX file which uses the SANDreport class file.
% It shows how a SAND report should be formatted, what sections and
% elements it should contain, and how to use the SANDreport class.
%
% Get the latest version of the class file and more at
%    https://gitlab.sandia.gov/latex/sandreport-latex
%
\makeatletter
\def\input@path{{../sand_report_template}}
\makeatother

% Bonus: 
% Scientific Notation Command (built-in)
%  This command allows for a standard way to adding units and scientific notation
%    eliminating the need for the constant use of \mathrm and explicitly typing 10^n 
%    every time you need a number in scientific notation 
%
%  The prototype command is:
%     \sci{mantissa}<exponent>[units]
%  Units will be typeset in mathrm so use standard notation (e.g. [g/cm^3])
%  other forms include:
%     * \sci{mantissa}<exponent>
%     * \sci{mantissa}[units]
%     * \sci{mantissa}
%******************* WARNING *********************
%
% Do not be concerned that there are a lot of 
% pages marked as 'intentionally left blank'.
% This is required by common-look-and-feel. 
%
% DO NOT attempt to change the class to remove 
% them! It will not pass common-look-and-feel review 
% without the blank pages marked and without 
% secctions/chapters starting on odd-numbered 
% pages.
%
% The class will continue to use the `article`
% style template as default, however, 
% `article`-style reports (where the top level 
% organizational unit is the `chapter`, 
% not `chapter`) will only be approved for 
% for SAND reports under 25 pages. 
%
% ************************************************
\documentclass[12pt,report,justified]{SANDreport}
\usepackage{minted}
\usepackage{xcolor}


% Define custom colors
\definecolor{codebg}{RGB}{245,245,245}      % light gray background
\definecolor{linegray}{RGB}{128,128,128}    % gray for line numbers
\definecolor{altlinebg}{RGB}{230,230,230}   % alternating line color

% Set global minted style
\setminted{
    %fontsize=\scriptsize,
    fontsize=\tiny,
    linenos,                        % show line numbers
    numbersep=8pt,                  % space between numbers and code
    frame=lines,                    % box frame
    framesep=2mm,
    bgcolor=codebg,                 % background color
    breaklines=true
}

% Style for line numbers
\renewcommand{\theFancyVerbLine}{%
    \textcolor{linegray}{\ttfamily\tiny\arabic{FancyVerbLine}}%
}




% ---------------------------------------------------------------------------- %
%
% Set the title, author, and date
%
    \title{A Trivially Short SAND Report (justified, report-style}

    \author{John E. Smith \& Jane S. Miller}

    % There is a "Printed" date on the title page of a SAND report, so
    % the generic \date should generally be empty.
    \date{}


% ---------------------------------------------------------------------------- %
% Set some things we need for SAND reports. These are mandatory
%
\SANDnum{SAND2002-xxxx}
\SANDprintDate{June 2002}

% ---------------------------------------------------------------------------- %
% The following definition does not have a default value and will not
% print anything, if not defined
%
\SANDsupersed{SAND1901-0001}{January 1901}


% ---------------------------------------------------------------------------- %
%
% Start the document
%
\begin{document}
    \maketitle

% ------------------------------------------------------------------------ %
% An Abstract is required for SAND reports
%
\begin{abstract}
  So abstract it hurts.
\end{abstract}


% ------------------------------------------------------------------------ %
% The table of contents and list of figures and tables
% Comment out \listoffigures and \listoftables if there are no
% figures or tables. Make sure this starts on an odd numbered page
%
\tableofcontents
\listoffigures
\listoftables


% ---------------------------------------------------------------------- %
% This is where the body of the report begins; usually with an Introduction
%
\SANDmain		% Start the main part of the report

\chapter{Introduction}
\label{Intro}
Very glad to meet you.


\section{A Section}
Underpromising, and hopefully over-delivering.

\section{Another Section}
Look at us go!

\chapter{Background}

\chapter{Objective}

\chapter{Data}

\section{C Frame Modal}

\subsection{\tt{cframe\_modal/jobs\_coarse.iga}}
\inputminted[stripnl=false]{json}{cframe_modal/jobs_coarse.json}

\subsection{\tt{cframe\_modal/jobs\_fine.iga}}
\inputminted[stripnl=false]{json}{cframe_modal/jobs_fine.json}

\subsection{\tt{cframe\_modal/jobs\_fine\_x2.iga}}
\inputminted[stripnl=false]{json}{cframe_modal/jobs_fine_x2.json}


\chapter{Methods}

\chapter{Results}

\chapter{Conclusion}
That is, as they say, all, folks.


\appendix
\chapter{Appendices}
    Include any supplementary material, such as additional data, calculations, or detailed explanations of methods that are too lengthy for the main sections.



\end{document}
